% !TeX spellcheck = es_CL
\documentclass[]{article}
\usepackage{geometry,tikz,lipsum,amsmath,amssymb,soul,xcolor,cancel,esint}
\tikzset{every picture/.style={line width=0.75pt}} %set default line width to 0.75pt        
\numberwithin{equation}{section}

\usepackage{../../../simonPackage}
%opening
\title{Ejercicios Certamen 3 Mecánica Clásica}
\author{Simon Fonseca}

\begin{document}
	\maketitle
	\section{Problema 1}
	The transformation of coordinates and momenta has a form
	\begin{align}
		P &= \alpha\frac{p^2 q}{\sqrt{1 + p^{2}q^{2}}}, & Q &= \beta \frac{\sqrt{1 + p^2 q^2}}{p}, & \alpha, \beta &= \mathrm{const}
	\end{align}
	\begin{enumerate}
		\item[a.] Find the values of constants $\alpha, \beta$ for which the transformation is canonical. For these values of $\alpha, \beta$, find explicitly the generating function $F$ which corresponds to the transformation.
		\resp Podemos calcular el corchete de Poisson entre $[P,Q]_{p,q}$, es cual deberá ser cero. Usando \texttt{Mathematica} tenemos
		\begin{equation}
			[P,Q]_{p,q} = \alpha\beta
		\end{equation}
		Por lo tanto, para que la transformación sea canónica deberá cumplirse que
		\begin{equation}
			\beta = \frac{1}{\alpha}
		\end{equation}
		
		\qed
		
		\item[b.] Apply the transformation found in the previous item to the hamiltonian of harmonic oscillator $H (p, q) = \frac{p^2}{2m} + \frac{kq^2}{2}$ and find the hamiltonian $\tilde{H} (P, Q)$ in terms of the new variables $(P, Q)$. Write out the canonical equations in terms of these new variables.
		\resp
		
		\qed
		
	\end{enumerate}
	
	
	
	
\end{document}